\makeatletter
\@ifundefined{keepfig}{\newenvironment{keepfig}{\par\edef\keepfig@pi{\the\parindent}\noindent
  \begin{minipage}{\linewidth}\parindent=\keepfig@pi\relax}{\end{minipage}\par}}{}
\makeatother
\section{Когда $G_f\subset A_n$?}

\subsection*{Знак перестановки и дискриминант}

Пусть $\sigma$~--- перестановка множества $\{1,2,\dots,n\}$. Пары $(i,j)$ с $i<j$, но $\sigma(i)>\sigma(j)$, называются
\emph{инверсиями} перестановки $\sigma$; перестановка \emph{чётна} или \emph{нечётна} в зависимости от чётности числа её
инверсий, а её \emph{знак} $\sgn(\sigma)$ равен $+1$ или $-1$ соответственно. Отображение $\sgn\colon S_n\to\{\pm1\}$~---
единственный гомоморфизм, переводящий каждую транспозицию в $-1$ (группа $S_n$ порождена транспозициями). Поэтому знак
перестановки $\sigma$ произвольного $n$-элементного множества $S$ можно определить, занумеровав $S$ и отождествив
$\sigma$ с перестановкой множества $\{1,\dots,n\}$: другая нумерация заменяет эту перестановку сопряжённой, а
гомоморфизм в абелеву группу $\{\pm1\}$ на сопряжённых элементах принимает одно и то же значение.

Пусть
\[
f(X)=X^n+a_1X^{n-1}+\dots+a_n\in F[X]
\]
\textup{(}$n\ge1$\textup{)}, и $f(X)=\prod_{i=1}^n(X-\alpha_i)$ в некотором поле разложения $F_f$. Положим
\[
\Delta(f)=\prod_{1\le i<j\le n}(\alpha_i-\alpha_j),\qquad D(f)=\Delta(f)^2=\prod_{1\le i<j\le n}(\alpha_i-\alpha_j)^2 .
\]
Элемент $D(f)$ называется \emph{дискриминантом} многочлена $f$. Очевидно, $D(f)\ne0$ тогда и только тогда, когда все
корни $f$ простые, т.\,е.\ $f$ сепарабелен. Для сепарабельного $f$ группа Галуа $G_f=\Gal(F_f/F)$ отождествляется с
подгруппой группы $\Sym(\{\alpha_1,\dots,\alpha_n\})$ (\rf{3.27a}), и для $\sigma\in G_f$ через $\sgn(\sigma)$ обозначается знак
соответствующей перестановки корней.

\begin{proposition}[4.1]
Пусть $f\in F[X]$~--- сепарабельный многочлен и $\sigma\in G_f$. Тогда
\begin{enumerate}[label=\textup{(\alph*)}]
\item $\sigma\Delta(f)=\sgn(\sigma)\,\Delta(f)$;
\item $\sigma D(f)=D(f)$.
\end{enumerate}
\end{proposition}

\begin{proof}
(a) Пусть $\sigma\alpha_i=\alpha_{\pi(i)}$, $\pi\in S_n$. Так как $\sigma$~--- гомоморфизм колец,
\[
\sigma\Delta(f)=\prod_{i<j}\bigl(\alpha_{\pi(i)}-\alpha_{\pi(j)}\bigr).
\]
Отображение $\{i,j\}\mapsto\{\pi(i),\pi(j)\}$~--- биекция множества двухэлементных подмножеств $\{1,\dots,n\}$ на себя. Поэтому
для каждой пары $k<l$ разность $\alpha_k-\alpha_l$ встречается справа ровно один раз~--- со знаком <<$+$>>, если она получилась
из пары $i<j$ с $\pi(i)<\pi(j)$, и со знаком <<$-$>>, если $(i,j)$~--- инверсия $\pi$. Значит,
$\sigma\Delta(f)=(-1)^{\#\{\text{инверсий }\pi\}}\Delta(f)=\sgn(\pi)\Delta(f)$.

(b) Возведём (a) в квадрат.
\end{proof}

Рассуждение использует лишь то, что $\sigma$~--- гомоморфизм колец, переставляющий корни. Иначе говоря, в
$\Z[X_1,\dots,X_n]$ для любой $\pi\in S_n$ выполнено тождество $\prod_{i<j}(X_{\pi(i)}-X_{\pi(j)})=\sgn(\pi)\prod_{i<j}(X_i-X_j)$.

\begin{remark}[4.1a]
(a) \emph{Зависимость от нумерации.} По тому же тождеству перенумерация корней перестановкой $\tau$ умножает $\Delta(f)$ на
$\sgn(\tau)$. Поэтому $\Delta(f)$ определён лишь с точностью до знака, а $D(f)$ от нумерации не зависит.

(b) \emph{Произведение по всем упорядоченным парам.} Часто удобнее произведение по всем парам $i\ne j$:
\[
\prod_{i\ne j}(\alpha_i-\alpha_j)=(-1)^{n(n-1)/2}D(f).
\]
Действительно, каждая неупорядоченная пара $\{i,j\}$ даёт слева множитель
$(\alpha_i-\alpha_j)(\alpha_j-\alpha_i)=-(\alpha_i-\alpha_j)^2$, а справа~--- $(\alpha_i-\alpha_j)^2$, и всего таких пар
$\binom n2=\frac{n(n-1)}2$.
\end{remark}

\begin{corollary}[4.2]
Пусть $f\in F[X]$~--- сепарабельный многочлен степени $n$, $F_f$~--- его поле разложения и $G_f=\Gal(F_f/F)$.
\begin{enumerate}[label=\textup{(\alph*)}]
\item Дискриминант $D(f)$ лежит в $F$.
\item Пусть $\operatorname{char}F\ne2$. Подполе поля $F_f$, отвечающее подгруппе $A_n\cap G_f$, есть $F[\Delta(f)]$. Поэтому
\[
G_f\subset A_n\iff\Delta(f)\in F\iff D(f)\ \text{является квадратом в}\ F .
\]
\end{enumerate}
\end{corollary}

\begin{proof}
(a) Расширение $F_f/F$~--- расширение Галуа (\rf{3.10}), а $D(f)\in F_f$ неподвижен относительно $G_f$ по \rf{4.1}(b);
поэтому $D(f)\in F_f^{G_f}=F$.

(b) Так как корни $f$ простые, $\Delta(f)\ne0$, а так как $\operatorname{char}F\ne2$, то $-\Delta(f)\ne\Delta(f)$. Поэтому
равенство $\sigma\Delta(f)=\sgn(\sigma)\Delta(f)$ из \rf{4.1}(a) показывает, что $\sigma\in G_f$ оставляет $\Delta(f)$ на месте
тогда и только тогда, когда $\sigma\in A_n$. Значит, $\Gal(F_f/F[\Delta(f)])=G_f\cap A_n$, и по основной теореме (\rf{3.17})
при соответствии Галуа
\[
G_f\cap A_n\ \longleftrightarrow\ F[\Delta(f)] .
\]
Следовательно, $G_f\subset A_n\iff G_f\cap A_n=G_f\iff F[\Delta(f)]=F\iff\Delta(f)\in F$. Наконец, если $D(f)=c^2$ с
$c\in F$, то $\Delta(f)^2=c^2$ и $\Delta(f)=\pm c\in F$; обратное очевидно.
\end{proof}

\begin{center}
\begin{keepfig}
\centering
\begin{tikzpicture}[every node/.style={font=\small}]
\node[cold] (F) at (0,0) {$F$};
\node[cnew] (M) at (0,1.5) {$F[\Delta(f)]=F[\sqrt{D(f)}]$};
\node[cold] (E) at (0,3) {$F_f$};
\draw (F) -- node[right] {$G_f/(G_f\cap A_n)$, порядок $1$ или $2$} (M);
\draw (M) -- node[right] {$G_f\cap A_n$} (E);
\draw[decorate,decoration={brace,amplitude=5pt,mirror}] (-1.75,3) -- node[left=6pt] {$G_f$} (-1.75,0);
\end{tikzpicture}
\figcap{Башня $F\subset F[\Delta(f)]\subset F_f$ (при $\operatorname{char}F\ne2$). Подгруппа $G_f\cap A_n$ имеет в $G_f$
индекс $1$ или $2$, поэтому $[F[\Delta(f)]:F]\le2$; равенство $F[\Delta(f)]=F$ означает, что $G_f\subset A_n$.}
\end{keepfig}
\end{center}

Утверждение (a) верно и без предположения сепарабельности (тогда $D(f)=0$), и вообще не требует теории Галуа:
$\prod_{i<j}(X_i-X_j)^2$~--- симметрический многочлен с целыми коэффициентами, поэтому
$D(f)=\Delta_n(a_1,\dots,a_n)$ для универсального многочлена $\Delta_n\in\Z[A_1,\dots,A_n]$ (\rf{5.37}(c)). Пункт (b) при
$\operatorname{char}F=2$ неверен: там $-1=1$, и $\Delta(f)$ не различает чётные и нечётные перестановки (см.\ \rf{4.4}).

\begin{example}[4.2a]
\emph{Степень $2$.} Корни $X^2+bX+c$ равны $\frac{-b\pm\sqrt{b^2-4c}}2$ (при $\operatorname{char}F\ne2$), поэтому
\[
\Delta(X^2+bX+c)=\pm\sqrt{b^2-4c},\qquad D(X^2+bX+c)=b^2-4c
\]
(последнее равенство верно в любой характеристике: $(\alpha_1-\alpha_2)^2=(\alpha_1+\alpha_2)^2-4\alpha_1\alpha_2$).

\emph{Степень $3$.} Для многочлена $X^3+bX+c$
\[
D(X^3+bX+c)=-4b^3-27c^2
\]
(шесть выводов~--- в \rf{4.2e} и \rf{4.37}). Сдвигом переменной (<<выделением полного куба>>) к этому виду приводится
любой кубический многочлен, если $\operatorname{char}F\ne3$ (\rf{4.2b}); общая формула:
\[
D(X^3+aX^2+bX+c)=a^2b^2-4b^3-4a^3c-27c^2+18abc .
\]

\emph{Высшие степени.} Хотя для корней многочлена степени $n>4$ нет общей формулы в радикалах (\rf{5.40b}), для его
дискриминанта формула есть всегда~--- это многочлен $\Delta_n$ от коэффициентов. Но он быстро усложняется: для
$X^4+aX^3+bX^2+cX+d$ в нём $16$ членов (\rf{4.9a}(c)), для $X^5+aX^4+bX^3+cX^2+dX+e$~--- уже $59$. Системы компьютерной
алгебры вычисляют его мгновенно (например, \texttt{poldisc(X\^{}3+a*X\^{}2+b*X+c,X)} в PARI); вручную дискриминант
удобнее всего находить через результант (\rf{4.35}--\rf{4.39}).
\end{example}

\begin{remark}[4.2b: сдвиг и растяжение]
(a) Для $t\in F$ многочлен $f(X+t)$ имеет корни $\alpha_i-t$ с теми же попарными разностями, поэтому
$D(f(X+t))=D(f)$.

(b) Для $\lambda\in F^\times$ многочлен $g(X)=\lambda^{-n}f(\lambda X)$ унитарен и имеет корни $\alpha_i/\lambda$, поэтому
$D(g)=\lambda^{-n(n-1)}D(f)$; для кубического многочлена это множитель $\lambda^{-6}$.

(c) \emph{Общий кубический многочлен.} Если $\operatorname{char}F\ne3$, то подстановка $X=T-\frac a3$ даёт
\[
f\bigl(T-\tfrac a3\bigr)=T^3+pT+q,\qquad p=b-\frac{a^2}3,\quad q=\frac{2a^3}{27}-\frac{ab}3+c,
\]
и по (a) $D(f)=-4p^3-27q^2$; раскрыв скобки, получаем формулу из \rf{4.2a}. По \rf{4.2c} она верна и при
$\operatorname{char}F=3$.

(d) \emph{Неунитарные многочлены.} Для $f=a_0\prod_{i=1}^n(X-\alpha_i)$, $a_0\ne0$, дискриминант обычно определяют как
$a_0^{2n-2}\prod_{i<j}(\alpha_i-\alpha_j)^2=a_0^{2n-2}D(f/a_0)$; множитель $a_0^{2n-2}$~--- квадрат, так что на ответ на вопрос
<<квадрат ли дискриминант?>> он не влияет.
Для кубики $AX^3+BX^2+CX+E$, подставив в (c) $a=B/A$, $b=C/A$, $c=E/A$ и умножив на $A^4$, получаем
\[
B^2C^2-4AC^3-4B^3E-27A^2E^2+18ABCE .
\]
\end{remark}

\begin{remark}[4.2c: принцип специализации]
Вычисляя дискриминант, удобно делить на коэффициенты (как в \rf{4.2e}, \rf{4.37}, \rf{4.39}) или предполагать, что
характеристика равна нулю. Это законно в силу универсальности многочлена $\Delta_n$. Пусть, скажем, коэффициенты
многочлена $f$ являются многочленами с целыми коэффициентами от параметров $T_1,\dots,T_k$ (например,
$f=X^n+T_1X+T_2$), и над полем рациональных функций $\Q(T_1,\dots,T_k)$ доказано равенство $D(f)=P(T_1,\dots,T_k)$, где
$P\in\Z[T_1,\dots,T_k]$. Левая часть равна $\Delta_n(a_1(T),\dots,a_n(T))\in\Z[T_1,\dots,T_k]$ (\rf{5.37}(c)), так что это
равенство многочленов с целыми коэффициентами. Применив гомоморфизм $\Z[T_1,\dots,T_k]\to F$, $T_i\mapsto t_i$, получаем
$D(f(t_1,\dots,t_k))=P(t_1,\dots,t_k)$ для \emph{любого} поля $F$ и \emph{любых} $t_i\in F$. В частности, при выводе формулы
можно делить на любой ненулевой многочлен от параметров (например, считать $p\ne0$, $q\ne0$ или $n(n-1)a\ne0$): над
$\Q(T)$ он обратим. Так же рассуждают в \rf{5.50a}(a).
\end{remark}

\subsection*{Дискриминант через производную и через критические значения}

\begin{proposition}[4.2d]
Пусть $f\in F[X]$~--- унитарный многочлен степени $n$ и $f=\prod_{i=1}^n(X-\alpha_i)$ в поле разложения.
\begin{enumerate}[label=\textup{(\alph*)}]
\item $f'(\alpha_i)=\prod_{j\ne i}(\alpha_i-\alpha_j)$, и
\[
D(f)=(-1)^{n(n-1)/2}\prod_{i=1}^nf'(\alpha_i).
\]
\item Пусть $n\cdot1\ne0$ в $F$ и $f'=n\prod_{j=1}^{n-1}(X-c_j)$ в некотором расширении. Тогда
\[
D(f)=(-1)^{n(n-1)/2}\,n^n\prod_{j=1}^{n-1}f(c_j).
\]
\end{enumerate}
\end{proposition}

\begin{proof}
(a) По правилу Лейбница $f'(X)=\sum_{k=1}^n\prod_{j\ne k}(X-\alpha_j)$. При $X=\alpha_i$ все слагаемые с $k\ne i$
обращаются в нуль (в них есть множитель $X-\alpha_i$), и остаётся $f'(\alpha_i)=\prod_{j\ne i}(\alpha_i-\alpha_j)$. Перемножая
по $i$, получаем $\prod_{i\ne j}(\alpha_i-\alpha_j)$, и остаётся применить \rf{4.1a}(b). (Это же рассуждение используется в
доказательстве \rf{5.49}.)

(b) Подставим $f'(\alpha_i)=n\prod_j(\alpha_i-c_j)$ и поменяем порядок произведений:
\[
\prod_{i=1}^nf'(\alpha_i)=n^n\prod_{j=1}^{n-1}\prod_{i=1}^n(\alpha_i-c_j)=n^n\prod_{j=1}^{n-1}(-1)^nf(c_j)
=n^n\prod_{j=1}^{n-1}f(c_j),
\]
так как $\prod_i(\alpha_i-c)=(-1)^n\prod_i(c-\alpha_i)=(-1)^nf(c)$, а $(-1)^{n(n-1)}=1$. Остаётся применить (a).
\end{proof}

Над $\R$ точки $c_j$~--- критические точки $f$, а $f(c_j)$~--- \emph{критические значения}; формула (b) говорит, что
дискриминант с точностью до множителя есть произведение критических значений. В частности, $D(f)=0$ тогда и только
тогда, когда некоторое критическое значение равно нулю, т.\,е.\ $f$ и $f'$ имеют общий корень~--- кратный корень $f$.

\begin{example}[4.2e: пять выводов формулы $D(X^3+pX+q)=-4p^3-27q^2$]
Пусть $f=X^3+pX+q=(X-x_1)(X-x_2)(X-x_3)$. По формулам Виета
\[
e_1=x_1+x_2+x_3=0,\qquad e_2=x_1x_2+x_1x_3+x_2x_3=p,\qquad e_3=x_1x_2x_3=-q .
\]

(i) \emph{Через определитель Вандермонда.} Квадрат определителя Вандермонда $V=(x_j^{\,i})$ равен
$\det(VV^{\mathsf T})$~--- определителю матрицы степенных сумм $s_{i+j}$, $s_k=\sum_ix_i^k$; этот вывод проведён в решении
упражнения 1-4. (Так же для любого $n$: $D(f)=\det(s_{i+j})_{0\le i,j\le n-1}$.)

(ii) \emph{Через производную.} По \rf{4.2d}(a) при $n=3$ имеем $D=-\prod_i f'(x_i)=-\prod_i(3x_i^2+p)$. Раскрывая скобки,
\[
\prod_i(3x_i^2+p)=27(x_1x_2x_3)^2+9p\,(x_1^2x_2^2+x_1^2x_3^2+x_2^2x_3^2)+3p^2(x_1^2+x_2^2+x_3^2)+p^3 .
\]
Здесь $(x_1x_2x_3)^2=e_3^2=q^2$, $\sum x_i^2=e_1^2-2e_2=-2p$ и
\[
x_1^2x_2^2+x_1^2x_3^2+x_2^2x_3^2=e_2^2-2e_1e_3=p^2
\]
(возведите в квадрат $e_2$: удвоенные произведения равны $2x_1x_2x_3(x_1+x_2+x_3)=2e_1e_3$). Итак,
$D=-(27q^2+9p^3-6p^3+p^3)=-4p^3-27q^2$.

(iii) \emph{Через критические значения.} Пусть $\operatorname{char}F\ne3$. Тогда $f'=3X^2+p=3(X-c)(X+c)$, где $c^2=-\frac p3$, и
по \rf{4.2d}(b) $D=-27f(c)f(-c)$. Так как $c^3=-\frac p3c$,
\[
f(\pm c)=\pm c\,(c^2+p)+q=q\pm\tfrac{2p}3c,\qquad f(c)f(-c)=q^2-\tfrac{4p^2}9c^2=q^2+\tfrac{4p^3}{27},
\]
откуда $D=-27q^2-4p^3$. Вычисление чисто алгебраическое: число $c$ лежит в $F[\sqrt{-p/3}]$ и не обязано быть
вещественным. Случай $\operatorname{char}F=3$ покрывается принципом \rf{4.2c}.

(iv) \emph{Через произведения $x_if'(x_i)$.} Будем считать, что $2pq\ne0$ в $F$ (общий случай сводится к этому
принципом специализации \rf{4.2c}; впрочем, в любой характеристике формула уже доказана в (ii)). Из $x_i^3=-px_i-q$
получаем
\[
x_if'(x_i)=3x_i^3+px_i=-2px_i-3q=-2p\,(x_i-s),\qquad s=-\frac{3q}{2p}.
\]
Перемножим эти равенства, учитывая $\prod_ix_i=-q$, $D=-\prod_if'(x_i)$ (по \rf{4.2d}(a)) и $\prod_i(x_i-s)=-f(s)$:
\begin{align*}
qD&=-q\prod_if'(x_i)=\prod_ix_if'(x_i)=(-2p)^3\prod_i(x_i-s)=8p^3f(s)\\
&=8p^3\Bigl(-\frac{27q^3}{8p^3}-\frac{3q}2+q\Bigr)=-27q^3-4p^3q,
\end{align*}
откуда $D=-4p^3-27q^2$. По существу это вычисление результанта $\operatorname{Res}(f',f)$ из \rf{4.37}, только без
самого слова <<результант>>.

(v) \emph{Исключение одного корня.} Так как $x_3=-(x_1+x_2)$, имеем $x_1-x_3=2x_1+x_2$ и $x_2-x_3=x_1+2x_2$. Положим
$P=x_1x_2$. Формулы Виета дают $p=P-(x_1+x_2)^2$ и $q=P(x_1+x_2)$, поэтому
\[
(x_1-x_2)^2=(x_1+x_2)^2-4P=-p-3P,\qquad (2x_1+x_2)(x_1+2x_2)=2(x_1+x_2)^2+P=3P-2p .
\]
Следовательно,
\[
D=(x_1-x_2)^2\bigl((2x_1+x_2)(x_1+2x_2)\bigr)^2=-(p+3P)(3P-2p)^2=-4p^3-27P^2(P-p)=-4p^3-27q^2,
\]
так как $P^2(P-p)=P^2(x_1+x_2)^2=q^2$. Этот вывод обходится без делений и верен в любой характеристике.

Шестой вывод~--- через результант~--- приведён в \rf{4.37}.
\end{example}

\begin{example}[4.2f: дискриминанты многочленов четвёртой степени]
(a) \emph{Биквадратный многочлен} $f=X^4+aX^2+b$. Его корни $\pm u$, $\pm v$, где $u^2+v^2=-a$, $u^2v^2=b$. Пары $\{u,-u\}$ и
$\{v,-v\}$ дают множители $(2u)^2(2v)^2=16b$, а четыре <<смешанные>> пары~--- $\bigl((u-v)(u+v)(-u-v)(-u+v)\bigr)^2=(u^2-v^2)^4$.
Так как $(u^2-v^2)^2=(u^2+v^2)^2-4u^2v^2=a^2-4b$,
\[
D(X^4+aX^2+b)=16\,b\,(a^2-4b)^2 .
\]

(b) \emph{Трёхчлен} $X^4+aX+b$: по \rf{4.39} при $n=4$
\[
D(X^4+aX+b)=256b^3-27a^4 .
\]

(c) Явные формулы для приведённой квартики $X^4+pX^2+qX+r$ и для общей $X^4+aX^3+bX^2+cX+d$ (в последней
$16$ членов) проще всего получить из кубической резольвенты, дискриминант которой равен дискриминанту квартики
(\rf{4.9}); они выписаны в \rf{4.9a}(c).
\end{example}

\subsection*{Знак дискриминанта над $\R$}

\begin{remark}[4.3]
Пусть $F\subset\R$ и $f\in F[X]$~--- сепарабельный многочлен степени $n$, у которого $n-2s$ вещественных корней и $s$
пар комплексно сопряжённых невещественных корней. Тогда
\[
\sgn D(f)=(-1)^s .
\]
Действительно, сгруппируем множители $(\alpha_i-\alpha_j)^2$ произведения $D(f)$. Пара вещественных корней даёт
положительный множитель. Пара сопряжённых корней $z,\bar z$ даёт $(z-\bar z)^2=(2i\operatorname{Im}z)^2=-4(\operatorname{Im}z)^2<0$.
Вещественный корень $r$ и пара $z,\bar z$ дают $(r-z)^2(r-\bar z)^2=\bigl((r-z)\overline{(r-z)}\bigr)^2=|r-z|^4>0$. Две
различные пары $\{z,\bar z\}$, $\{w,\bar w\}$ дают
$(z-w)^2(\bar z-\bar w)^2\cdot(z-\bar w)^2(\bar z-w)^2=|z-w|^4|z-\bar w|^4>0$. Отрицательных множителей ровно $s$.

Отрицательное число не является квадратом в $F\subset\R$, поэтому \emph{при нечётном $s$ группа $G_f$ не содержится в
$A_n$} (\rf{4.2}(b)). Это можно увидеть и прямо: \emph{комплексное сопряжение действует на корнях как произведение $s$
независимых транспозиций}, т.\,е.\ является нечётной перестановкой, и гомоморфизм $\sgn\colon G_f\to\{\pm1\}$ сюръективен.
Почему комплексное сопряжение $c$ вообще лежит в $G_f$? Возьмём в качестве поля разложения подполе
$F_f=F[\alpha_1,\dots,\alpha_n]\subset\CC$, где $\alpha_i\in\CC$~--- корни $f$ (поле разложения единственно с точностью до
$F$-изоморфизма, \rf{2.13}). Автоморфизм $c$ поля $\CC$ тождествен на $\R\supset F$, поэтому $f(\bar\alpha)=\overline{f(\alpha)}=0$
для каждого корня $\alpha$: $c$ переставляет корни $f$, а значит, $c(F_f)=F[\bar\alpha_1,\dots,\bar\alpha_n]=F_f$. Ограничение
$c|_{F_f}$~--- автоморфизм $F_f$ над $F$, т.\,е.\ элемент $G_f$; он оставляет на месте вещественные корни и меняет местами
корни каждой невещественной пары. (Условие $F\subset\R$ существенно: при $F=\Q[i]$ сопряжение не тождественно на $F$.)

Обратное неверно: при чётном $s$ группа $G_f$ не обязана лежать в $A_n$. Например, $X^3-4X+1$ неприводим над $\Q$ (корни
$\pm1$ не подходят), имеет три вещественных корня ($s=0$), но $D=-4(-4)^3-27=229$ не квадрат, так что $G_f=S_3$. У
$X^4+2$ все четыре корня невещественны ($s=2$), $D=256\cdot2^3=2^{11}$ (\rf{4.2f}(b)) не квадрат, и $G_f\not\subset A_4$.
\end{remark}

\begin{example}[4.3a: вещественная кубика]
Пусть $f=X^3+pX+q$, $p,q\in\R$, $D=-4p^3-27q^2$. Тогда
\begin{itemize}[leftmargin=*]
\item $D>0$ $\iff$ у $f$ три различных вещественных корня;
\item $D=0$ $\iff$ у $f$ кратный корень; тогда все корни вещественны;
\item $D<0$ $\iff$ у $f$ один вещественный и два комплексно сопряжённых невещественных корня.
\end{itemize}
При $D\ne0$ это частный случай \rf{4.3} ($s\in\{0,1\}$). Если $D=0$, то кратный корень вещественен: вместе с невещественным
корнем кратности $\ge2$ корнем той же кратности был бы и сопряжённый, а всего корней с учётом кратности три; тогда и
оставшийся корень вещественен.

Формула \rf{4.2e}(iii) делает это наглядным. При $p<0$ критические точки $\mp c$, $c=\sqrt{-p/3}>0$, вещественны: в $-c$
локальный максимум, в $c$~--- локальный минимум, и график трижды пересекает ось тогда и только тогда, когда
$f(-c)>0>f(c)$, т.\,е.\ $f(c)f(-c)<0$, т.\,е.\ $D=-27f(c)f(-c)>0$. При $p\ge0$ функция $f$ монотонна ($f'=3x^2+p\ge0$), у неё
ровно один вещественный корень, и $D=-4p^3-27q^2\le0$ (равенство~--- лишь при $p=q=0$, тройной корень $0$).
\end{example}

\begin{center}
\begin{keepfig}
\centering
\begin{tikzpicture}[every node/.style={font=\small}]
\begin{scope}[xscale=1.1,yscale=0.55]
\draw[->,gray] (-2.6,0) -- (2.7,0) node[right] {$x$};
\draw[->,gray] (0,-3.4) -- (0,5.6);
\draw[cold,thick,domain=-2.15:2.25,samples=90] plot (\x,{\x*\x*\x-3*\x+0.6});
\draw[cnew,dashed] (-1,0) -- (-1,2.6); \draw[cnew,dashed] (1,0) -- (1,-1.4);
\fill[cnew] (-1,2.6) circle (2pt) node[above] {$f(-c)>0$};
\fill[cnew] (1,-1.4) circle (2pt) node[below right] {$f(c)<0$};
\node[cnew,below] at (-1,0) {$-c$}; \node[cnew,above] at (1,0) {$c$};
\fill (0,0.6) circle (1.5pt) node[above right=-1pt] {$q$};
\foreach \r in {-1.8245,0.2028,1.6217}{\fill[cold] (\r,0) circle (1.6pt);}
\node[cold,below left] at (-1.82,0) {$x_1$};
\node[cold,above right=-1pt] at (0.2,0) {$x_2$};
\node[cold,below right] at (1.62,0) {$x_3$};
\end{scope}
\begin{scope}[xshift=7.6cm,scale=0.9]
\draw[->,gray] (-1.8,0) -- (2.6,0) node[right] {$\operatorname{Re}$};
\draw[->,gray] (0,-2.3) -- (0,2.4) node[above] {$\operatorname{Im}$};
\coordinate (r) at (1.6,0); \coordinate (z) at (-0.8,1.5); \coordinate (zb) at (-0.8,-1.5);
\draw[cnew,->] (r) -- (z) node[midway,above right=-2pt] {$z-r$};
\draw[cnew,->] (r) -- (zb) node[midway,below right=-2pt] {$\bar z-r$};
\draw[cnew,dashed] (z) -- (zb) node[pos=0.25,left] {$z-\bar z$};
\fill[cold] (r) circle (2pt) node[below right] {$r$};
\fill[cold] (z) circle (2pt) node[above left] {$z$};
\fill[cold] (zb) circle (2pt) node[below left] {$\bar z$};
\end{scope}
\end{tikzpicture}
\figcap{Слева: кубика $f=x^3+px+q$ с $p<0$ имеет три вещественных корня, когда её критические значения $f(-c)$ и $f(c)$
разных знаков, т.\,е.\ $D=-27f(c)f(-c)>0$; точка $(0,q)$~--- точка перегиба, центр симметрии графика. Справа: при одном
вещественном корне $r$ множитель $(z-\bar z)^2$ отрицателен ($z-\bar z$ чисто мнимо), а $(r-z)^2(r-\bar z)^2=|r-z|^4>0$,
поэтому $D<0$.}
\end{keepfig}
\end{center}

\subsection*{Характеристика $2$}

\begin{aside}[4.4]
Если $\operatorname{char}F=2$, то дискриминант \emph{всегда} квадрат в $F$, и для вопроса <<лежит ли $G_f$ в $A_n$?>> он
бесполезен. Действительно, теперь $-1=1$, так что
\[
\Delta(f)=\prod_{i<j}(\alpha_i-\alpha_j)=\prod_{i<j}(\alpha_i+\alpha_j),
\]
а это выражение симметрично по корням. Поэтому $\Delta(f)\in F$: для сепарабельного $f$~--- по \rf{4.1}(a) (теперь
$\sgn(\sigma)=\pm1=1$ в $F$, и $\Delta(f)$ неподвижен относительно всей $G_f$), а в общем случае~--- без теории Галуа, так как
$\prod_{i<j}(X_i+X_j)$~--- симметрический многочлен с целыми коэффициентами (\rf{5.37}(c)). Например, для $X^2+bX+c$ имеем
$\Delta=\alpha_1+\alpha_2=b$, а для $X^3+aX^2+bX+c$
\[
\Delta=(\alpha_1+\alpha_2)(\alpha_1+\alpha_3)(\alpha_2+\alpha_3)=e_1e_2-e_3=ab+c,
\]
и действительно, по модулю $2$ формула из \rf{4.2a} даёт $a^2b^2+c^2=(ab+c)^2$. Так, $X^2+X+1$ неприводим над $\F_2$, его
группа Галуа $S_2\not\subset A_2$, а $D=1$~--- квадрат (общий квадратный случай~--- в \rf{4.6}(b)).

Вместо $\Delta(f)$ нужно взять элемент, на котором нечётные перестановки действуют <<сдвигом на $1$>>, а не сменой знака.
Для сепарабельного $f$ положим $t_{ij}=\alpha_i/(\alpha_i+\alpha_j)$ при $i\ne j$ (знаменатель $\alpha_i+\alpha_j=\alpha_i-\alpha_j$
не равен нулю) и $\beta=\sum_{i<j}t_{ij}$. Так как $t_{ij}+t_{ji}=1$, элемент $\sigma\in G_f$, переставляющий корни по
перестановке $\pi$, переводит слагаемое $t_{ij}$ в слагаемое $\beta$ для пары $\{\pi i,\pi j\}$, если $\pi i<\pi j$, и в это
слагаемое плюс $1$, если $(i,j)$~--- инверсия $\pi$. Поэтому $\sigma\beta=\beta$ для чётных и $\sigma\beta=\beta+1$ для нечётных
$\sigma$. Далее, $t_{ij}^2+t_{ij}=t_{ij}(t_{ij}+1)=t_{ij}t_{ji}$, и, так как в характеристике $2$ возведение в квадрат
аддитивно,
\[
\beta^2+\beta=\sum_{i<j}t_{ij}t_{ji}=D_{\mathrm B}(f)\overset{\mathrm{def}}{=}\sum_{i<j}\frac{\alpha_i\alpha_j}{\alpha_i^2+\alpha_j^2}
\]
(\emph{дискриминант Берлекэмпа}). Так как $(\beta+1)^2+(\beta+1)=\beta^2+\beta$, элемент $D_{\mathrm B}(f)$ неподвижен
относительно $G_f$, т.\,е.\ лежит в $F$. Отсюда
\[
G_f\subset A_n\iff\beta\in F\iff\text{существует }\delta\in F\text{ с }\delta^2+\delta=D_{\mathrm B}(f):
\]
первая равносильность верна, так как $F_f^{G_f}=F$; во второй импликация <<$\Rightarrow$>> очевидна ($\delta=\beta$), а
если $\delta^2+\delta=\beta^2+\beta$, то $(\delta+\beta)(\delta+\beta+1)=0$, т.\,е.\ $\beta\in\{\delta,\delta+1\}\subset F$
(то же рассуждение~--- в \rf{5.32e}). Вместо
извлечения квадратного корня (теория Куммера) здесь решается уравнение Артина--Шрайера $X^2+X=D_{\mathrm B}$
(ср.\ \rf{5.32c}). Для квадратного многочлена $X^2+X+a$ с корнями $\alpha$, $\alpha+1$ получается
$D_{\mathrm B}=\alpha(\alpha+1)/\bigl(\alpha^2+(\alpha+1)^2\bigr)=a$; так, для $X^2+X+1$ над $\F_2$ имеем $D_{\mathrm B}=1$, а
$\delta^2+\delta=1$ не имеет решений в $\F_2$, что и отвечает $G_f\not\subset A_2$.
\end{aside}

\section{Когда $G_f$ действует на корнях транзитивно?}

\begin{proposition}[4.5]
Пусть $f\in F[X]$~--- сепарабельный многочлен. Тогда $f$ неприводим в том и только в том случае, когда $G_f$ действует на
корнях $f$ транзитивно, т.\,е.\ для любых двух корней $\alpha,\beta$ найдётся $\sigma\in G_f$ с $\sigma\alpha=\beta$.
\end{proposition}

\begin{proof}
Это частный случай \rf{3.27d}: орбиты $G_f$ на множестве корней~--- в точности множества корней неприводимых
множителей $f$. Напомним идею. Если $g$~--- неприводимый множитель $f$ и $\alpha$~--- его корень, то $\sigma\alpha$ для
$\sigma\in G_f$ снова корень $g$ (так как $g(\sigma\alpha)=\sigma(g(\alpha))=0$), поэтому при транзитивном действии все корни $f$
являются корнями $g$, и $f=g$ (с точностью до множителя из $F$). Обратно, если $f$ неприводим и $\alpha,\beta$~--- его корни, то
$F[\alpha]$ и $F[\beta]$~--- оба поля корня $f$ (\rf{1.27a}), и существует $F$-изоморфизм $\varphi_0\colon F[\alpha]\to F[\beta]$,
$\alpha\mapsto\beta$ (\rf{2.1}); так как $F_f$~--- поле разложения $f$ и над $F[\alpha]$, и над $F[\beta]$, он продолжается до
$\sigma\in G_f$ (\rf{2.13a}; в доказательстве \rf{2.12a}, на котором оно основано, гомоморфизм продолжается шаг за
шагом, по одному корню, с помощью \rf{2.4}(b)~--- см.\ рисунок там):
\[
\begin{tikzcd}[row sep=small]
F_f\arrow[r,dashed,cnew,"\sigma"] & F_f\\
F[\alpha]\arrow[u,dash]\arrow[r,cold,"\varphi_0"',"\alpha\mapsto\beta"] & F[\beta]\arrow[u,dash]\\
 F\arrow[u,dash]\arrow[r,equal] & F\arrow[u,dash]
\end{tikzcd}
\]
\end{proof}

Если $f$ неприводим степени $n$, то $n=[F[\alpha]:F]$ делит $[F_f:F]=(G_f:1)$. Итак, $G_f$~--- \emph{транзитивная подгруппа
$S_n$, порядок которой делится на $n$}; дальше в этой главе именно среди таких подгрупп мы и будем искать группу Галуа.

% Приложение о результантах (4.35--4.39) --- в конце главы, файл s6_overQ.tex.
